\section{Related Work}
\subsection{KV Cache Eviction Algorithm}
% 在long-context的场景下，KV cache会随着输入长度的增加而增加，这会导致巨大的时间和空间开销。一些已有的方法通过维护一个固定大小的KV cache来保存一些关键的tokens，并驱逐剩余不重要的tokens来减少内存使用且减少计算量。

To save the significant time and space overhead as the input length increases, various approaches explore evicting tokens to reduce both memory usage and computational cost. StreamingLLM~\cite{streamingllm} introduces the phenomenon of "attention sink" and supports longer sequence by retaining only the KV pairs of attention sinks and recent tokens. H2O~\cite{h2o} employs a low-cost eviction strategy to maintain a fixed-size KV cache containing heavy-hitters, based on the sum of historical attention scores. SnapKV~\cite{snapkv} uses the last tokens in the prompt during the prefilling stage to select critical tokens for generation in the decoding stage. PyramidKV~\cite{pyramidkv}, PyramidInfer~\cite{pyramidinfer}, and LazyLLM~\cite{lazyllm} leverage attention distribution across layers to dynamically adjust cache size, 
% allocating more to lower layers and less to higher layers, thereby 
making cache allocation more efficient. 
Other methods like MagicPig~\cite{magicpig}, Q-Hitter~\cite{q-hitter}, ALISA~\cite{alisa} and , which combine KV cache eviction with quantization, hashing algorithms, or sparse windows, can also improve inference efficiency.
% However, all of these methods lead to varying degrees of performance degradation, as they permanently evict certain tokens that may play a crucial role in future generation processes.
However, these methods induce non-negligible performance degradation, as they potentially evict certain tokens that are crucial for future generation.

\subsection{KV Cache Retrieval Algorithm}
Unlike KV cache eviction algorithms, KV cache retrieval algorithms retain a complete KV cache and dynamically retrieve important KV pairs to reduce inference latency. 
% Quest~\cite{quest} uses PageAttention~\cite{vllm}, dividing tokens into pages, and applying an approximate attention score to retrieve the most relevant pages for the current decoding steps. 
Following PageAttention~\cite{vllm}, Quest~\cite{quest} divides tokens into pages and devises an approximate attention score to retrieve the most relevant pages for the current decoding steps.
InfLLM~\cite{infllm} adopts a strategy similar to Quest, offloading most of the cache to the CPU to support longer prompts.
% ShadowKV~\cite{shadowkv} enhances the storage efficiency of Quest by storing only low-rank key caches and offloading value caches, reducing memory usage and enabling inference with large batch sizes and long contexts. 
ShadowKV~\cite{shadowkv} enhances the storage efficiency of Quest by storing only low-rank key caches and offloading value caches, allowing inference with large batch sizes and context lengths.
Others like SparQ~\cite{sparq}, InfiniGen~\cite{infinigen}, and Loki~\cite{loki}, accelerate the selection of top-\textit{k} critical tokens by reducing the dimension. 
KV cache retrieval methods maintain performance by preserving the entire KV cache but inevitably incur increased inference latency and storage costs.

\subsection{Head-wise Attention Classification}
% Another line of work focuses on classifying attention heads based on their characteristics, with the goal of designing more efficient algorithms to reduce inference latency. 
Another line of work classifies attention heads into distinct sparse patterns.
MInference~\cite{minference} divides the attention into A-shape, Vertical-Slash, and Block-Sparse patterns, achieving acceleration during the \emph{pre-filling} stage. RazorAttention~\cite{razorattention} and DuoAttention~\cite{duoattention} split heads into retrieval heads and streaming heads to determine whether to implement full attention or an attention mechanism similar to StreamingLLM~\cite{streamingllm}. Methods like AdaKV~\cite{adakv} and HeadKV~\cite{headkv} achieve more fine-grained classification by allocating different budgets to each attention head.
These methods focus on the features within individual heads or their variations, but overlook the patterns of attention heads across the prefilling and decoding stages.